% AI-written, AI cross-reviewed; no human sign-off; not a claim that RH is proved
% RH Notes blueprint source excerpts for Z1.5.02
% Verbatim excerpts in the notation of the parent blueprint. The labelled editions and source ranges appear below.
% See the unit YAML for the current proof status. Earlier-round passages retain their original wording.

% Start of atlas edition annotation. The following block is outside the original excerpt.
\paragraph{Status of named inputs in this reading.}
\noindent\nolinkurl{GlobalTheory.T.iwasawaDecomposition_finiteAdeles_GL} --- in-repo constant (\path{Development/AbelianVarietyTheory/Proofs/IwasawaDecomposition/IwasawaDecompositionFiniteAdelesGLStatement.lean:14--15}); not compared.\par
\noindent\nolinkurl{UpperTheory.V.Af.clozelCohC_alt} --- in-repo constant (\path{Development/AutomorphicArithmetic/AutomorphicVarieties/ModularityArgument/ThresholdArguments/DenseThresholdCases/TPCohC.lean:10--30}); not compared.\par
External theorem numbers in the original excerpt are not verified here unless a separate edition-specific source record documents the comparison.
% End of atlas edition annotation.

\section*{Z1.5.02. Counting automorphic eigensystems}

% Parent source SHA-256: 1af64922b98ed2310744411b1f18b25cc426c97220646321471836e1a3a409d6
% blueprint/ch_P.tex lines 12522 through 12531; r5:small:Z1.5.02
\subsection{Z1.5.02. Counting a finite family at principal level}
\label{r5:small:Z1.5.02}
\textit{Complete with the stabilizer and coinduction argument on page~\pageref{r6:Z1.5.02}.} Fix a CM field $K$ and $n\ge2$. A finite family $S\subset\mathrm{SigK}(K,m,M)$ consists of cusp representations with the prescribed infinitesimal character at each embedding, conjugate duality and a nonzero vector fixed by $K(M)$. Assume isomorphism implies equality in $S$, $M\ge1$, and $m$ strictly decreasing with entries bounded by $H$ in absolute value. Take the single coefficient representation $\Xi$ just constructed, with $\dim\Xi\le c_1(H+1)^{A_1}$. The ideal $(3M)$ lies in $(M)$, so the reduction modulo $M$ factors through that modulo $3M$ and $K(3M)\subseteq K(M)$. Every given fixed vector survives at that smaller subgroup. The cohomological comparison applied to each member with this $\Xi$, followed by the injectivity bound for this one finite family, bounds $|S|$ by the same total cohomology dimension at $3M$.
\par{\footnotesize \nolinkurl{UpperTheory.V.Af.SigK} (\nolinkurl{Development/AutomorphicArithmetic/AutomorphicVarieties/AutomorphicForms/IsoOfFrob.lean}\allowbreak:13--14). \nolinkurl{UpperTheory.V.Af.CardK} (\nolinkurl{Development/AutomorphicArithmetic/AutomorphicVarieties/AutomorphicForms/IsoOfFrob.lean}\allowbreak:19--23). \nolinkurl{UpperTheory.V.Af.card_k_of_printed_results} (\nolinkurl{Development/AutomorphicArithmetic/AutomorphicVarieties/ModularityArgument/ThresholdArguments/DenseThresholdCases/TCardK.lean}\allowbreak:18--36). \nolinkurl{klev_three} (\nolinkurl{Development/AutomorphicArithmetic/AutomorphicVarieties/ModularityArgument/ThresholdArguments/DenseThresholdCases/TCardK.lean}\allowbreak:10--17).\par}

To see where the exponent in the cohomology estimate comes from, use the map from the finite-adele cosets at $3M$ to those at $3$. Every fiber has the same size $[K(3):K(3M)]$. It is bounded by $(3M)^{[K:\mathbb Q]n^2}$ by embedding the congruence quotient at level one into the set of matrices over the finite quotient of the integral adeles. The base has finitely many orbits for the full rational group; the stabilizer over $gK(3)$ is the rational intersection with $gK(3)g^{-1}$. Decompose the coefficient functions over these orbits. On a fiber the dimension of the coefficient representation grows by at most the fiber size, by taking a copy at every point. The cohomology bound for its stabilizer, followed by the orbit/coinduction comparison, then bounds the total by $c_2 M^{A_2}\dim\Xi$, with $A_2=[K:\mathbb Q]n^2$. As $M,H+1\ge1$, increasing both exponents to $A_1+A_2$ yields $|S|\le c_1c_2[M(H+1)]^{A_1+A_2}$.
\par{\footnotesize \nolinkurl{gen} (\nolinkurl{Development/AutomorphicArithmetic/AutomorphicVarieties/AutomorphicForms/TemperedLocalComponents/CohBoundFree.lean}\allowbreak:28--60). \nolinkurl{cohBound_of} (\nolinkurl{Development/AutomorphicArithmetic/AutomorphicVarieties/AutomorphicForms/TemperedLocalComponents/CohBoundFree.lean}\allowbreak:69--91). \nolinkurl{UpperTheory.V.Af.cohBoundFree} (\nolinkurl{Development/AutomorphicArithmetic/AutomorphicVarieties/AutomorphicForms/TemperedLocalComponents/CohBoundFree.lean}\allowbreak:93--99). \nolinkurl{UpperTheory.level_subgroup_index_bound_of_card_ohat} (\nolinkurl{Development/AutomorphicArithmetic/AutomorphicVarieties/BuildingsResidueChars/LevelSubgroupIndexBoundOfCardOhat.lean}\allowbreak:22--38). \nolinkurl{UpperTheory.V.Af.card_k_of_printed_results} (\nolinkurl{Development/AutomorphicArithmetic/AutomorphicVarieties/ModularityArgument/ThresholdArguments/DenseThresholdCases/TCardK.lean}\allowbreak:30--36).\par}

\textit{Stabilizer comparison, proved on page~\pageref{r6:Z1.5.02}.} The comparison concerns the arithmetic stabilizer bound $\exists C\ \forall V,\ \mathrm{totCoh}(\Gamma,V)\le C\operatorname{rank}_{\mathbb C}V$ for groups commensurable with the image of $\mathrm{GL}_n(\mathcal O_K)$, and its coinduction comparison bounding the sum over orbits. These are project dependencies of the cited cohomology bound. \textit{[Hypothesis, discharged at the construction on page~\pageref{r7:review:Z1.5.02}.]} The cohomological comparison for cusp representations and the injection for pairwise nonisomorphic members are respectively the fields \nolinkurl{cohC} and \nolinkurl{borelC} of the input record. Z1.5.03 passes this record unchanged when it invokes the count. The constructor supplies both fields as described on page~\pageref{r7:review:Z1.5.02}.
\par{\footnotesize \nolinkurl{UpperTheory.V.Af.bddArith} (\nolinkurl{Development/AutomorphicArithmetic/AutomorphicVarieties/AutomorphicForms/TemperedLocalComponents/BddArith.lean}\allowbreak:59--72). \nolinkurl{UpperTheory.V.Af.CohB.Bdd} (\nolinkurl{Development/AutomorphicArithmetic/AutomorphicVarieties/AutomorphicForms/HomologyRank/RankHomology.lean}\allowbreak:17). \nolinkurl{UpperTheory.V.Af.cohBoundFree} (\nolinkurl{Development/AutomorphicArithmetic/AutomorphicVarieties/AutomorphicForms/TemperedLocalComponents/CohBoundFree.lean}\allowbreak:93--99). \nolinkurl{UpperTheory.V.Af.Hypotheses'} (\nolinkurl{Development/AutomorphicArithmetic/AutomorphicVarieties/ModularityArgument/ThresholdArguments/DenseThresholdCases/AutomorphicFormHyp.lean}\allowbreak:12--22). \nolinkurl{UpperTheory.GalN.V.Win.N.countK_alt} (\nolinkurl{Development/AutomorphicArithmetic/AutomorphicVarieties/ModularityArgument/ThresholdArguments/DenseThresholdCases/FrobeniusTraceCountBoundOfCmEigensystemCount.lean}\allowbreak:63--125).\par}

% blueprint/ch_P.tex lines 12663 through 12698; r6:Z1.5.02
\subsection{Z1.5.02. Why the arithmetic cohomology bound is uniform}\label{r6:Z1.5.02}
Fix the CM field $K$ and $n\ge2$ of \nolinkurl{UpperTheory.V.Af.CardK}. For a complex representation $V$ of a discrete group $H$ put
\[
 b(H,V)=\sum_{q\ge0}\dim_{\mathbf C}H^q(H,V).
\]
The sum and ranks here are cardinals, so the statement also makes sense for infinite-dimensional $V$. We prove $b(H,V)\le C_H\dim V$ for the stabilizers needed at level three. The constants may depend on the finitely many base orbits. They are chosen before the finer level $3M$, its coefficient module and the family of representations. This is the totally complex branch of \nolinkurl{UpperTheory.V.Af.bddArith}; it uses the same Hermitian complex as Proposition~\ref{prop:P-r4-k7}. {\footnotesize\nolinkurl{UpperTheory.V.Af.CohB.totCoh, Bdd}, \nolinkurl{Development/AutomorphicArithmetic/AutomorphicVarieties/AutomorphicForms/HomologyRank/RankHomology.lean}\allowbreak:14--17.} {\footnotesize\nolinkurl{UpperTheory.V.Af.bddArith}, \nolinkurl{Development/AutomorphicArithmetic/AutomorphicVarieties/AutomorphicForms/TemperedLocalComponents/BddArith.lean}\allowbreak:59--72.}

\paragraph{Uniformly many nonzero degrees.}
Let $G=\mathrm{GL}_n(\mathcal O_K)$ and let $X_q$ be the ordered $(q+1)$-tuples of vertices whose normal cones meet in the positive cone. Faces omit one vertex; repetitions are allowed. K7 gives exactness of the augmented integral tuple complex in every degree, finitely many $G$-orbits at each degree and finite vertex stabilizers. Tensor with $\mathbf C$; exactness persists by flatness over $\mathbf Z$. Each tuple has finite stabilizer $S$ because $S$ fixes its first vertex. Each orbit summand $\mathbf C[G/S]$ is projective: a lift over a surjection is specified by a lift of the first basis vector, which one averages over $S$ to make it $S$-fixed. Thus $C_q=\mathbf C[X_q]$ gives a projective resolution of the trivial module. A $G$-map $C_q\to V$ is determined on a finite set of orbit representatives, so
$\dim H^q(G,V)\le |G\backslash X_q|\dim V$. A bound on the number of contributing degrees is still essential.
{\footnotesize\nolinkurl{UpperTheory.bounded_complex_of_bounded_simplicial_complex}, \nolinkurl{Development/AutomorphicArithmetic/AutomorphicVarieties/BuildingsResidueChars/BoundedComplexOfBoundedSimplicialComplex.lean}\allowbreak:8--28.} {\footnotesize\nolinkurl{UpperTheory.permutation_module_homology_rank}, \nolinkurl{Development/AutomorphicArithmetic/AutomorphicVarieties/BuildingsResidueChars/PermutationModuleHomologyRank.lean}\allowbreak:4--34.} {\footnotesize\nolinkurl{UpperTheory.V.Af.bddArith_cx}, \nolinkurl{Development/AutomorphicArithmetic/AutomorphicVarieties/AutomorphicForms/TemperedLocalComponents/BddArith.lean}\allowbreak:38--49.}

Choose a finite set $R$ of ordered edge orbit representatives, including the degenerate edge orbit whenever needed. Fix the first vertex $v$ of a tuple. For each $r\in R$ whose first vertex is in the orbit of $v$, choose $k_r$ carrying $v$ to that vertex $v_r$. If the tuple has distinct vertices $v_j$, carry $(v,v_j)$ to its representative $r_j$ by $g_j$. The pair $(r_j,g_jk_{r_j}^{-1})$ lies in $\coprod_{r\in R}G_{v_r}$ and recovers the edge: apply $(g_jk_{r_j}^{-1}k_{r_j})^{-1}$ to $r_j$. This injects the $q+1$ indices into a finite set of size $D=\sum_r |G_{v_r}|$. Every tuple of degree $q\ge D$ therefore repeats a vertex.
{\footnotesize\nolinkurl{UpperTheory.adelic_rep_dimension_bounded}, \nolinkurl{Development/AutomorphicArithmetic/AutomorphicVarieties/BuildingsResidueChars/AdelicRepDimensionBounded.lean}\allowbreak:4--68.}

Apply the normalized signed average of all permutations of $q+1$ indices to $C_q$. This commutes with the differential: on deleting index $i$ of a permuted tuple, the image index $j$ and the induced permutation of the remaining indices have sign $(-1)^{i+j}$ times the old sign. The $q+1$ possible choices account for the change in the factorial in the denominator. At degree zero this chain map is the identity. Above the cutoff the transposition of repeated entries cancels each term with another term of opposite sign. A chain map of projective resolutions lifting the identity is homotopic to the identity: subtract it from the identity, lift the resulting map into the kernel of the augmentation through the first differential, and continue by exactness and projectivity in each degree. Consequently the zero map at $q\ge D$ also induces the identity on $H^q(G,V)$, which forces this cohomology to vanish. Sum the orbit counts for $q<D$ to obtain $C_G$. In particular the tuple complex itself need not vanish in high degrees.
{\footnotesize\nolinkurl{UpperTheory.rep_hom_alternating_sum_eq_zero}, \nolinkurl{Development/AutomorphicArithmetic/AutomorphicVarieties/BuildingsResidueChars/RepHomAlternatingSumEqZero.lean}\allowbreak:10--58.} {\footnotesize\nolinkurl{UpperTheory.altChain_of_permSplit}, \nolinkurl{Development/AutomorphicArithmetic/AutomorphicVarieties/BuildingsResidueChars/AltChainOfPermSplit.lean}\allowbreak:5--46.} {\footnotesize\nolinkurl{UpperTheory.ext_vanishing_of_resolution_endo}, \nolinkurl{Development/AutomorphicArithmetic/AutomorphicVarieties/BuildingsResidueChars/ExtVanishingOfResolutionEndo.lean}\allowbreak:32--38.} {\footnotesize\nolinkurl{UpperTheory.adelic_rep_bounded_of_degree_bounded}, \nolinkurl{Development/AutomorphicArithmetic/AutomorphicVarieties/BuildingsResidueChars/AdelicRepBoundedOfDegreeBounded.lean}\allowbreak:16--26.}

\paragraph{Moving the bound to an arithmetic stabilizer.}
For $S\le G$ of finite index $h$, a coinduced function $F$ with $F(sg)=sF(g)$ is determined by values on $h$ coset representatives. Hence $\dim\operatorname{Coind}_S^G V=h\dim V$. Degreewise Shapiro follows by restricting a free resolution of the trivial $G$-module to $S$ (the group algebra is free over $\mathbf C[S]$) and evaluating the coinduction adjunction on each free term. Thus $C_S=hC_G$ suffices. In the opposite direction an $H$-module $W$, for $S\le H$ of finite index $d$, is a summand of its restricted and then coinduced module. The inclusion is $w\mapsto(g\mapsto gw)$. A left inverse sends $F$ to
$d^{-1}\sum_{gS\in H/S}gF(g^{-1})$. Changing $g$ to $gs$ leaves each summand unchanged; left translation permutes the cosets and proves equivariance. Shapiro gives $b(H,W)\le b(S,W|_S)$. Applying these two inequalities to the intersection transfers the bound to every subgroup commensurable with the image of $G$.
{\footnotesize\nolinkurl{UpperTheory.adelic_rep_bounded_of_subgroup_bounded}, \nolinkurl{Development/AutomorphicArithmetic/AutomorphicVarieties/BuildingsResidueChars/AdelicRepBoundedOfSubgroupBounded.lean}\allowbreak:4--57.} {\footnotesize\nolinkurl{UpperTheory.coind_retract_of_finite_index}, \nolinkurl{Development/AutomorphicArithmetic/AutomorphicVarieties/BuildingsResidueChars/CoindRetractOfFiniteIndex.lean}\allowbreak:5--75.} {\footnotesize\nolinkurl{UpperTheory.adelic_rep_bounded_over_of_sum_bound_coind_retract}, \nolinkurl{Development/AutomorphicArithmetic/AutomorphicVarieties/BuildingsResidueChars/AdelicRepBoundedOverOfSumBoundCoindRetract.lean}\allowbreak:4--27.} {\footnotesize\nolinkurl{UpperTheory.adelic_rep_bounded_comm_of_bounded_equiv}, \nolinkurl{Development/AutomorphicArithmetic/AutomorphicVarieties/BuildingsResidueChars/AdelicRepBoundedCommOfBoundedEquiv.lean}\allowbreak:3--17.}

For the counting application take $\Gamma=\mathrm{GL}_n(K)$ acting through the rational inclusion on $Y=\mathrm{GL}_n(\mathbb A_{K,f})/K(3)$. The stabilizer of $gK(3)$ is the inverse image of $gK(3)g^{-1}$. Two compact open subgroups are commensurable because the open cosets of their intersection cover each compact subgroup. Inverse image preserves finite relative indices; the preimage of $K(1)$ consists of matrices whose entries and inverse entries are integral at every finite place, exactly $\mathrm{GL}_n(\mathcal O_K)$. There are finitely many base orbits. At level one the upper triangular decomposition and approximation for its unipotent factor reduce to diagonal matrices; for each diagonal entry choose an idele ideal-class representative, absorb its field element on the left and its integral unit on the right. Each level-one double coset splits into finitely many level-three double cosets, since $[K(1):K(3)]$ is finite.
{\footnotesize\nolinkurl{UpperTheory.arithmetic_subgroups_commensurable}, \nolinkurl{Development/AutomorphicArithmetic/AutomorphicVarieties/BuildingsResidueChars/ArithmeticSubgroupsCommensurable.lean}\allowbreak:4--48.} {\footnotesize\nolinkurl{UpperTheory.arithmetic_subgroup_level_one}, \nolinkurl{Development/AutomorphicArithmetic/AutomorphicVarieties/BuildingsResidueChars/ArithmeticSubgroupLevelOne.lean}\allowbreak:6--41.} {\footnotesize\nolinkurl{PIW.ClassFinD.borelKGL_of}, \nolinkurl{Formalism/PeriodIntegrals/UnitaryGroupMassFormula/IdeleClassFiniteness/BorelGL.lean}\allowbreak:83--102.} {\footnotesize\nolinkurl{UpperTheory.finite_adelic_double_cosets}, \nolinkurl{Development/AutomorphicArithmetic/AutomorphicVarieties/BuildingsResidueChars/FiniteAdelicDoubleCosets.lean}\allowbreak:13--21.} {\footnotesize\nolinkurl{UpperTheory.adelic_cohomology_finite_orbit_bound}, \nolinkurl{Development/AutomorphicArithmetic/AutomorphicVarieties/BuildingsResidueChars/AdelicCohomologyFiniteOrbitBound.lean}\allowbreak:3--19.}

\paragraph{The same finite fiber in the coinduction comparison.}
Put $X=\mathrm{GL}_n(\mathbb A_{K,f})/K(3M)$ and let $p:X\to Y$ be projection. Above one base orbit of $y$, write $X_y=p^{-1}(y)$, $S=\Gamma_y$, and $V=\mathbf C^d$ with action $\xi$. For functions on the inverse image of the orbit the required isomorphism to $\operatorname{Coind}_S^\Gamma\operatorname{Map}(X_y,V)$ is
\[
 (Tf)(g)(z)=\xi(g) f(g^{-1}z).
\]
It satisfies $(Tf)(sg)=s(Tf)(g)$ and transforms $\Gamma$-translation of $f$ into right translation in $g$. Given $F$ in the target, choose $g$ with $gp(x)=y$ and set $f(x)=\xi(g)^{-1}F(g)(gx)$. Any other choice is $sg$ with $s\in S$ and the displayed equivariance cancels $s$. Substitution gives both inverse identities. The fiber has $h=[K(3):K(3M)]$ points. Its function space therefore has dimension $hd$; Shapiro and the stabilizer bound give at most $C_Shd$ in total cohomology. There are finitely many base orbits, so the total is at most $hd\sum_{[y]} C_{\Gamma_y}$. All these stabilizers belong to the \emph{fixed} level.
{\footnotesize\nolinkurl{UpperTheory.V.Af.CohB.toCo}, \nolinkurl{Development/AutomorphicArithmetic/AutomorphicVarieties/AutomorphicForms/HomologyRank/RankHomology.lean}\allowbreak:33--60.} {\footnotesize\nolinkurl{UpperTheory.cohomology_to_coinvariants_surjective}, \nolinkurl{Development/AutomorphicArithmetic/AutomorphicVarieties/BuildingsResidueChars/CohomologyToCoinvariantsSurjective.lean}\allowbreak:4--32.} {\footnotesize\nolinkurl{UpperTheory.shapiro_functorial_of_coinduction_bijective}, \nolinkurl{Development/AutomorphicArithmetic/AutomorphicVarieties/BuildingsResidueChars/ShapiroFunctorialOfCoinductionBijective.lean}\allowbreak:4--21.} {\footnotesize\nolinkurl{UpperTheory.absolute_bound_of_partial_bound_of_orbit_bound}, \nolinkurl{Development/AutomorphicArithmetic/AutomorphicVarieties/BuildingsResidueChars/AbsoluteBoundOfPartialBoundOfOrbitBound.lean}\allowbreak:4--34.}

Finally $h\le(3M)^{[K:\mathbf Q]n^2}$ by reducing matrices modulo $3M$; absorbing the fixed power of three and using the common coefficient representation of Z1.5.01 gives the polynomial count in Section~\ref{r5:small:Z1.5.02}. \textit{Complete under the public assumptions.} The fields \nolinkurl{cohC} and \nolinkurl{borelC} are hypotheses of the counting theorem: the former supplies the required cohomological contribution for the chosen weight, and the latter bounds the cardinality of the same family by the cohomology dimension. The stabilizer and fiber estimates above are deductions internal to that application.
{\footnotesize\nolinkurl{UpperTheory.V.Af.cohBoundFree}, \nolinkurl{Development/AutomorphicArithmetic/AutomorphicVarieties/AutomorphicForms/TemperedLocalComponents/CohBoundFree.lean}\allowbreak:69--99.} {\footnotesize\nolinkurl{UpperTheory.V.Af.card_k_of_printed_results}, \nolinkurl{Development/AutomorphicArithmetic/AutomorphicVarieties/ModularityArgument/ThresholdArguments/DenseThresholdCases/TCardK.lean}\allowbreak:18--36.}

% Parent source SHA-256: 1fdc6ea9aac6cf10316f85b3253b33ae98289a76ff5aa2e79daaaf11e057ce18
% blueprint/ch_P_r7.tex lines 3 through 10; r7:review:Z1.5.02
\subsection{The finite class set used in Z1.5.02}\label{r7:review:Z1.5.02}
The count on page~\pageref{r6:Z1.5.02} uses a constant for the finitely many arithmetic stabilizers already present at level three. Here is the finiteness argument at level one, followed by the passage to the open compact subgroups used in the proof. Put $G=\mathrm{GL}_n$, $U$ its upper unitriangular group and $K_0=G(\widehat{\mathcal O}_K)$. At a finite place, right integral column operations send the last row of an invertible matrix to $(0,\ldots,0,a)$. Choose an entry of minimum valuation as pivot, move it to the last position, and subtract integral multiples of this column from the rest. Repeat on the upper left invertible block. This expresses the matrix as an upper triangular matrix times an integral invertible matrix. At places where the matrix already belongs to $K_0$, take the triangular factor to be the identity. Hence $g=utk$, with $u\in U(\mathbb A_f)$, $t$ diagonal and $k\in K_0$. \nolinkurl{GlobalTheory.T.iwasawaDecomposition_finiteAdeles_GL} (\path{Development/AbelianVarietyTheory/Proofs/IwasawaDecomposition/IwasawaDecompositionFiniteAdelesGLStatement.lean:14--15}).

The strictly upper entries give coordinates on $U(\mathbb A_f)$ and its rational points. In each coordinate $K$ is dense in the finite adeles. Indeed, clear the finitely many local denominators by a positive integer $d$, require congruences to sufficiently high powers at the specified primes, and impose divisibility by $d$ at additional primes dividing $d$. The Chinese remainder theorem produces an integral numerator satisfying all these congruences. Its quotient by $d$ is integral elsewhere. The open neighborhood $u(tK_0t^{-1}\cap U(\mathbb A_f))$ therefore contains some $\nu\in U(K)$. Then $ut=\nu tw$ with $w\in K_0$. \nolinkurl{GlobalTheory.T.strongApproximation_upperUnitriangular_finiteAdeles} (\path{Development/AbelianVarietyTheory/Proofs/IwasawaDecomposition/StrongApproximationUpperUnitriangularFiniteAdeleStatement.lean:14--15}).

The valuations of a diagonal entry give a fractional ideal. Their simultaneous kernel is $\widehat{\mathcal O}_K^\times$, and principal fractional ideals are the images of $K^\times$. The ideal class group is finite \cite[Theorem 4.4]{P-R3-ANT}. Choose one idele for each ideal class by prescribing its finitely supported valuations. Splitting each diagonal entry into a rational scalar, a chosen idele and an integral unit leaves at most $h_K^n$ diagonal representatives. This proves finiteness of $G(K)\backslash G(\mathbb A_f)/K_0$. Passing to level three subdivides each representative into at most $[K_0:K(3)]$ classes, finite by reduction modulo three. The three inputs and their assembly are \nolinkurl{PIW.ClassFinD.idele_class_finite} (\path{Formalism/PeriodIntegrals/UnitaryGroupMassFormula/IdeleClassFiniteness/LatticeIdealNonzero.lean:43--108}), \nolinkurl{PIW.ClassFinD.borelKGL_of} (\path{Formalism/PeriodIntegrals/UnitaryGroupMassFormula/IdeleClassFiniteness/BorelGL.lean:83--102}) and \nolinkurl{UpperTheory.finite_adelic_double_cosets} (\path{Development/AutomorphicArithmetic/AutomorphicVarieties/BuildingsResidueChars/FiniteAdelicDoubleCosets.lean:13--21}).

\emph{[Hypothesis, discharged at the automorphic-record construction.]} In \nolinkurl{UpperTheory.V.Af.card_k_of_printed_results} (\path{Development/AutomorphicArithmetic/AutomorphicVarieties/ModularityArgument/ThresholdArguments/DenseThresholdCases/TCardK.lean:18--36}), $\mathrm{cohC}$ identifies the cohomological contribution for the specified coefficient, and $\mathrm{borelC}$ injects the family into the dimension estimate at level $3M$. This principal congruence subgroup is open and compact: it is a kernel of reduction to a finite discrete matrix group inside the compact open integral group, with $3M\ge1$ (\nolinkurl{UpperTheory.V.Af.isOpen_KLev}, \nolinkurl{UpperTheory.V.Af.isCompact_KLev} (\path{Development/AutomorphicArithmetic/AutomorphicVarieties/AutomorphicForms/KLevTop.lean:66--91})). The direct caller \nolinkurl{UpperTheory.GalN.V.Win.N.countK_alt} (\path{Development/AutomorphicArithmetic/AutomorphicVarieties/ModularityArgument/ThresholdArguments/DenseThresholdCases/FrobeniusTraceCountBoundOfCmEigensystemCount.lean:63--125}) passes the record through. The construction \nolinkurl{UpperTheory.V.HV.DescentVanishingReductions.published_hypotheses_from_patching_alt} (\path{Development/AutomorphicArithmetic/AutomorphicVarieties/ModularityArgument/ThresholdArguments/DenseThresholdCases/PublishedHypothesesFromPatching.lean:98--103}) builds the extended record; its projection \nolinkurl{UpperTheory.V.Af.HypothesesH'.toClassicalTheorems} (\path{Development/AutomorphicArithmetic/AutomorphicVarieties/ModularityArgument/ThresholdArguments/DenseThresholdCases/AfDesc.lean:17--23}) gives the required record. Its cohomological field comes from the explicit literature-rank-$N$ premise via \nolinkurl{UpperTheory.V.Af.clozelCohC_alt} (\path{Development/AutomorphicArithmetic/AutomorphicVarieties/ModularityArgument/ThresholdArguments/DenseThresholdCases/TPCohC.lean:10--30}). The other field is \nolinkurl{UpperTheory.V.Z.Coh.borelInjC} (\path{Development/AutomorphicArithmetic/AutomorphicVarieties/ModularityArgument/ClozelCohomology/ClozelConjC.lean:31}). Together with the bound on page~\pageref{r6:Z1.5.02} this yields $c_2c_1[M(H+1)]^{A_1+A_2}$, with $A_2=[K:\mathbf Q]n^2$.

